\documentclass[tikz,border=6mm]{standalone}
\usepackage[UTF8]{ctex}
\usepackage{xcolor}
\usetikzlibrary{positioning,fit,backgrounds,calc}

% ---------- Color palette ----------
\definecolor{TitleBlue}{HTML}{1F4E79}
\definecolor{PanelBorder}{HTML}{B7C8D8}
\definecolor{PanelBg}{HTML}{FBFCFE}
\definecolor{TextDark}{HTML}{263746}
\definecolor{TextMuted}{HTML}{5F6B76}

\definecolor{BlueFill}{HTML}{E8F1FB}
\definecolor{BlueBorder}{HTML}{2F6DA3}
\definecolor{PurpleFill}{HTML}{F2EAFE}
\definecolor{PurpleBorder}{HTML}{7A4FB3}
\definecolor{OrangeFill}{HTML}{FFF2E2}
\definecolor{OrangeBorder}{HTML}{C97A1A}
\definecolor{RedFill}{HTML}{FDEBEC}
\definecolor{RedBorder}{HTML}{C84C61}
\definecolor{GreenFill}{HTML}{EAF7EA}
\definecolor{GreenBorder}{HTML}{4B8B3B}

% ---------- Content ----------
\newcommand{\ChartTitle}{Linux 内核驱动 Contract / Invariant 设计检查表}
\newcommand{\ChartSubtitle}{写代码之前，先回答这 10 个问题}

\newcommand{\QOneTitle}{1. Ownership / 所有权}
\newcommand{\QOneMain}{谁拥有这个资源？}
\newcommand{\QOneHint}{谁负责创建、引用、转移与最终释放？}

\newcommand{\QTwoTitle}{2. Lifetime / 生命周期}
\newcommand{\QTwoMain}{这个资源什么时候有效？}
\newcommand{\QTwoHint}{从 probe 到 remove/reset，创建与销毁边界在哪里？}

\newcommand{\QThreeTitle}{3. Concurrency / 并发}
\newcommand{\QThreeMain}{谁能同时访问它？}
\newcommand{\QThreeHint}{process / IRQ / worker / timer 是否可能并发访问同一状态？}

\newcommand{\QFourTitle}{4. Locking / 锁规则}
\newcommand{\QFourMain}{访问时必须持什么锁？}
\newcommand{\QFourHint}{锁保护哪些字段？锁顺序是什么？是否存在 ABBA 风险？}

\newcommand{\QFiveTitle}{5. Context / 执行上下文}
\newcommand{\QFiveMain}{代码运行在哪种上下文？}
\newcommand{\QFiveHint}{process / softirq / hardirq / workqueue / timer？}

\newcommand{\QSixTitle}{6. Blocking / 阻塞规则}
\newcommand{\QSixMain}{这里允许 sleep 吗？}
\newcommand{\QSixHint}{能否 mutex\_lock、等待 completion、分配 GFP\_KERNEL 内存？}

\newcommand{\QSevenTitle}{7. Error Path / 失败路径}
\newcommand{\QSevenMain}{失败以后必须恢复到什么状态？}
\newcommand{\QSevenHint}{部分初始化如何回滚？资源释放是否完整且顺序正确？}

\newcommand{\QEightTitle}{8. Teardown / 拆卸顺序}
\newcommand{\QEightMain}{remove / reset 时按什么顺序停止？}
\newcommand{\QEightHint}{先阻止新请求，再停 IRQ / DMA / worker，最后释放资源。}

\newcommand{\QNineTitle}{9. Hardware State / 硬件状态}
\newcommand{\QNineMain}{软件状态如何对应硬件状态？}
\newcommand{\QNineHint}{寄存器、DMA、IRQ、BAR 与 driver state 是否保持一致？}

\newcommand{\QTenTitle}{10. Invariants / 不变量}
\newcommand{\QTenMain}{哪些关系绝不能被破坏？}
\newcommand{\QTenHint}{例如：IRQ 仍可能运行时，device context 绝不能被释放。}

% ---------- Styles ----------
\tikzset{
  titlebox/.style={
    rounded corners=3mm,
    draw=TitleBlue,
    fill=TitleBlue!5,
    line width=0.9pt,
    align=center,
    inner xsep=8mm,
    inner ysep=4mm,
    text width=158mm
  },
  card/.style={
    rounded corners=2.6mm,
    draw,
    line width=0.85pt,
    minimum width=77mm,
    minimum height=28mm,
    text width=68mm,
    align=left,
    inner xsep=5mm,
    inner ysep=4mm,
    anchor=north west
  },
  bluecard/.style={card, fill=BlueFill, draw=BlueBorder},
  purplecard/.style={card, fill=PurpleFill, draw=PurpleBorder},
  orangecard/.style={card, fill=OrangeFill, draw=OrangeBorder},
  redcard/.style={card, fill=RedFill, draw=RedBorder},
  greencard/.style={card, fill=GreenFill, draw=GreenBorder},
  footer/.style={
    rounded corners=2.5mm,
    draw=PanelBorder,
    fill=PanelBg,
    line width=0.8pt,
    align=center,
    text width=158mm,
    inner xsep=6mm,
    inner ysep=4mm
  }
}

\newcommand{\CardText}[3]{%
  {\bfseries\color{TextDark}#1}\\[1.2mm]
  {\large\bfseries #2}\\[1mm]
  {\footnotesize\color{TextMuted}#3}%
}

\begin{document}
\begin{tikzpicture}[node distance=6mm and 7mm]

% ---------- Title ----------
\node[titlebox] (title) {
  {\Large\bfseries\color{TitleBlue}\ChartTitle}\\[1mm]
  {\normalsize\color{TitleBlue!85!black}\ChartSubtitle}
};

% ---------- Row 1 ----------
\node[bluecard, below left=8mm and 0mm of title.south west, anchor=north west] (q1) {
  \CardText{\QOneTitle}{\QOneMain}{\QOneHint}
};
\node[bluecard, right=7mm of q1.north east, anchor=north west] (q2) {
  \CardText{\QTwoTitle}{\QTwoMain}{\QTwoHint}
};

% ---------- Row 2 ----------
\node[purplecard, below=6mm of q1.south west, anchor=north west] (q3) {
  \CardText{\QThreeTitle}{\QThreeMain}{\QThreeHint}
};
\node[purplecard, below=6mm of q2.south west, anchor=north west] (q4) {
  \CardText{\QFourTitle}{\QFourMain}{\QFourHint}
};

% ---------- Row 3 ----------
\node[orangecard, below=6mm of q3.south west, anchor=north west] (q5) {
  \CardText{\QFiveTitle}{\QFiveMain}{\QFiveHint}
};
\node[orangecard, below=6mm of q4.south west, anchor=north west] (q6) {
  \CardText{\QSixTitle}{\QSixMain}{\QSixHint}
};

% ---------- Row 4 ----------
\node[redcard, below=6mm of q5.south west, anchor=north west] (q7) {
  \CardText{\QSevenTitle}{\QSevenMain}{\QSevenHint}
};
\node[redcard, below=6mm of q6.south west, anchor=north west] (q8) {
  \CardText{\QEightTitle}{\QEightMain}{\QEightHint}
};

% ---------- Row 5 ----------
\node[greencard, below=6mm of q7.south west, anchor=north west] (q9) {
  \CardText{\QNineTitle}{\QNineMain}{\QNineHint}
};
\node[greencard, below=6mm of q8.south west, anchor=north west] (q10) {
  \CardText{\QTenTitle}{\QTenMain}{\QTenHint}
};

% ---------- Footer ----------
\node[footer, below=8mm of $(q9.south)!0.5!(q10.south)$] (footer) {
  {\bfseries\color{TitleBlue}使用原则：}\\[1mm]
  Contract 定义“什么操作在什么条件下合法”；Invariant 定义“哪些关系在整个生命周期中绝不能被破坏”。\\[1mm]
  {\footnotesize\color{TextMuted}建议在 probe/remove、IRQ、DMA、reset、hotplug、error path 编码前逐项回答，并把答案写进设计文档、注释和测试计划。}
};

% ---------- Background ----------
\begin{scope}[on background layer]
\node[
  fit=(q1)(q2)(q3)(q4)(q5)(q6)(q7)(q8)(q9)(q10)(footer),
  rounded corners=4mm,
  draw=PanelBorder,
  fill=white,
  line width=0.8pt,
  inner xsep=7mm,
  inner ysep=7mm
] {};
\end{scope}

\end{tikzpicture}
\end{document}
